\documentclass[11pt]{article}
\usepackage[margin=1in]{geometry}
\usepackage[T1]{fontenc}
\usepackage[utf8]{inputenc}
\usepackage{amsmath,amssymb,amsthm}
\usepackage{enumitem}

\title{ICPC World Finals 2020\\I. Quests}
\author{}
\date{}

\begin{document}
\maketitle

\section*{Problem Summary}

Each quest has base XP $x$ and target level $d$.
If we complete it while our current level is below $d$, we earn $cx$ instead of $x$.
We may do the quests in any order, and we want to maximize the final total XP.

Since the base $x$ is earned no matter what, the real problem is to maximize the total base-XP value of
the quests that still receive the multiplier.

\section*{Key Observations}

\begin{itemize}[leftmargin=*]
    \item If a quest receives the multiplier, all quests done before it also receive the multiplier.
    Therefore in an optimal order, all multiplied quests come first, and all non-multiplied quests come
    after them.
    \item Let $Y$ be the sum of the base XP values of the multiplied quests already completed.
    Then the actual XP earned so far is exactly $cY$.
    So quest $i$ with parameters $(x_i,d_i)$ can still be multiplied exactly when
    \[
        cY < d_i v.
    \]
    \item This is equivalent to
    \[
        Y \le \left\lfloor \frac{d_i v - 1}{c} \right\rfloor.
    \]
    So if we only look at multiplied quests, quest $i$ behaves like a job with:
    \begin{itemize}[leftmargin=*]
        \item processing time $x_i$,
        \item value $x_i$,
        \item latest allowed \emph{start time}
        \[
            A_i = \left\lfloor \frac{d_i v - 1}{c} \right\rfloor.
        \]
    \end{itemize}
    \item Turning a start-time bound into an ordinary completion deadline gives
    \[
        D_i = A_i + x_i.
    \]
    Then a set of multiplied quests is feasible iff it can be scheduled so that every selected job finishes
    by its deadline $D_i$.
    \item This is now a standard deadline knapsack:
    after sorting by deadline, if total selected processing becomes $s$, then quest $i$ may be the last
    selected quest exactly when
    \[
        s \le D_i.
    \]
\end{itemize}

\section*{Algorithm}

\begin{enumerate}[leftmargin=*]
    \item For every quest compute
    \[
        D_i = \left\lfloor \frac{d_i v - 1}{c} \right\rfloor + x_i.
    \]
    Also compute the base total
    \[
        B = \sum x_i.
    \]
    \item Sort the quests by increasing $D_i$.
    \item Let \texttt{reachable[s]} mean:
    after processing some prefix of the sorted quests, there exists a feasible set of multiplied quests whose
    total base-XP sum is exactly $s$.
    Initially only \texttt{reachable[0]} is true.
    \item Process quests in sorted order. For quest $i$, update the bitset as a $0/1$ knapsack transition:
    from every reachable sum $s-x_i$, we may reach $s$ if
    \[
        s \le D_i.
    \]
    \item Let $S$ be the largest reachable sum after all quests. Then the final answer is
    \[
        B + (c-1)S,
    \]
    because every multiplied quest contributes an extra $(c-1)x_i$ above its base reward.
\end{enumerate}

\section*{Correctness Proof}

We prove that the algorithm returns the correct answer.

\paragraph{Lemma 1.}
There exists an optimal playthrough in which every multiplied quest is done before every non-multiplied
quest.

\paragraph{Proof.}
Every quest always contributes its base reward $x$.
Doing a non-multiplied quest earlier can only increase the current XP and level, which can never help a
later quest receive the multiplier.
So if a non-multiplied quest appears before a multiplied one, swapping them cannot decrease the reward of
either quest.
By repeatedly applying this swap, we obtain an optimal order with all multiplied quests first. \qed

\paragraph{Lemma 2.}
If quest $i$ is multiplied after previously multiplied quests with total base-XP sum $Y$, then this is
possible exactly when
\[
    Y \le \left\lfloor \frac{d_i v - 1}{c} \right\rfloor.
\]

\paragraph{Proof.}
By Lemma 1, only multiplied quests matter before quest $i$.
If their base-XP sum is $Y$, the actual XP already earned is $cY$.
Quest $i$ is multiplied exactly when the current level is less than $d_i$, i.e.
\[
    cY < d_i v.
\]
Since all quantities are integers, this is equivalent to
\[
    Y \le \left\lfloor \frac{d_i v - 1}{c} \right\rfloor.
\]
\qed

\paragraph{Lemma 3.}
A set of multiplied quests is feasible iff, after giving each quest deadline
\[
    D_i = \left\lfloor \frac{d_i v - 1}{c} \right\rfloor + x_i,
\]
it can be scheduled so that every selected quest finishes by its deadline.

\paragraph{Proof.}
For a selected quest $i$, let $Y$ be the total base-XP sum of earlier selected quests.
By Lemma 2, quest $i$ is feasible exactly when $Y \le A_i$, where
\[
    A_i = \left\lfloor \frac{d_i v - 1}{c} \right\rfloor.
\]
Since quest $i$ itself contributes $x_i$, this is equivalent to
\[
    Y + x_i \le A_i + x_i = D_i.
\]
But $Y+x_i$ is exactly the completion time of quest $i$ in the scaled scheduling model. \qed

\paragraph{Lemma 4.}
After sorting quests by nondecreasing deadlines, the bitset DP computes exactly all feasible total
base-XP sums of multiplied quests.

\paragraph{Proof.}
For ordinary single-machine scheduling with deadlines, every feasible subset is feasible in
nondecreasing-deadline order.
So when processing quests in that order, a total sum $s$ is feasible exactly when there exists a feasible
sum $s-x_i$ from earlier quests and
\[
    s \le D_i.
\]
This is exactly the transition implemented by the DP.
Therefore the final bitset contains precisely all feasible multiplied-base-XP totals. \qed

\paragraph{Theorem.}
The algorithm outputs the correct answer for every valid input.

\paragraph{Proof.}
By Lemma 1, we may restrict attention to a prefix of multiplied quests.
By Lemma 3, choosing that prefix is equivalent to choosing a feasible subset in the deadline model.
By Lemma 4, the DP finds the maximum possible total base-XP sum $S$ of such a subset.
Each quest always contributes its base reward, and each multiplied quest contributes an additional
$(c-1)x$.
Hence the maximum final XP is exactly
\[
    \sum x_i + (c-1)S,
\]
which is what the algorithm outputs. \qed

\section*{Complexity Analysis}

Let
\[
    X = \sum x_i \le 4\,000\,000.
\]
The bitset has $X+1$ states.
Each of the $n$ transitions is a bitset shift-and-OR, so the running time is
\[
    O\!\left(n \cdot \frac{X}{64}\right),
\]
and the memory usage is $O(X)$ bits.

\section*{Implementation Notes}

\begin{itemize}[leftmargin=*]
    \item The DP is done with a manual dynamic bitset for speed.
    \item The scaled deadline
    \[
        \left\lfloor \frac{d_i v - 1}{c} \right\rfloor
    \]
    must use 64-bit arithmetic.
    \item Sorting by deadline and then by smaller $x$ is harmless and keeps equal-deadline handling stable.
\end{itemize}

\end{document}
